% ===================================================================
% THEORY — Evaluative and Delegative Trust
% ===================================================================
% Converted from docs/theoretical_framework_moves_1_2.md
%              and docs/theoretical_framework_moves_3_6.md
% ===================================================================

% --- Move 1: The Unidimensional Problem (no subsection heading) ---

A foundational finding of modern financial economics is that trust matters for economic outcomes. \citet{GuisoSapienza2004, GuisoSapienza2008} demonstrate that generalized trust predicts stock market participation, portfolio allocation, and the use of formal financial instruments, even after controlling for risk aversion, wealth, and education. The mechanisms linking trust to financial behavior have since received sustained empirical attention: \citet{SapenzaZingales2012} distinguish between generalized social trust and specific beliefs about institutional trustworthiness; \citet{GennaioliShleiferVishny2015} model financial advisors as ``money doctors'' whose primary service is providing trust that reduces clients' perceived variance of returns; and \citet{GurunStoffmanYonker2018} demonstrate that geographic and cultural proximity --- channels plausibly related to trust --- predict advisory relationships and their outcomes.

Despite the richness of these empirical contributions, the literature shares a common structural assumption: trust is treated as unidimensional. It varies in \textit{degree} --- investors trust more or less --- but not in \textit{kind}. The standard measurement approaches reflect this assumption. The World Values Survey item used by \citeauthor{GuisoSapienza2004} asks whether ``most people can be trusted'' on a binary scale. The trust game paradigm introduced by \citet{BergDickhautMcCabe1995} measures trust as the dollar amount sent to an anonymous counterparty --- a continuous quantity representing the magnitude of risk accepted. \citeauthor{MayerDavisSchoorman1995}'s~(\citeyear{MayerDavisSchoorman1995}) influential organizational framework decomposes trust into ability, benevolence, and integrity, but treats each as a continuous attribute of the trustee, not as a qualitatively distinct mode of the trust relationship itself.

These measurement approaches share a deeper assumption that \citet{hardin1993street} identifies as a feature of ``street-level'' trust epistemology: initial trust is a learned Bayesian capacity that the trustor brings \textit{to} the encounter, relatively independent of the specific trustee's properties. Before we accumulate evidence about a particular advisor or institution, we rely on priors learned from a category of interaction --- prior experiences with financial professionals, narratives about algorithmic decision-making, cultural attitudes toward institutional authority. \citeauthor{hardin1993street}'s insight reframes the measurement problem. Instruments that present a novel counterparty and measure the trustor's response are not capturing a reaction to that counterparty's characteristics; they are measuring the trustor's prior disposition toward the category. The question is not merely whether these instruments capture enough \textit{variance} in trust --- it is whether they capture the right \textit{kind}. A trustor who approaches algorithms as tools to be assessed and a trustor who approaches them as agents to be entrusted with discretion are not expressing different magnitudes of the same attitude. They are operating in different modes, learned through different histories of encounter, and the standard instruments register both as a single quantity.

This unidimensional treatment creates an explanatory gap. Consider two well\-documented patterns in algorithmic finance. First, \citet{DietvorstSimmonsMassey2015} find that individuals who observe an algorithm err --- even when it outperforms human judgment on the same task --- subsequently prefer human forecasters. This ``algorithm aversion'' appears to reflect a fragile form of trust in which a single negative signal triggers exit. Second, and apparently in tension, \citet{LoggMinsonMoore2019} find that individuals systematically \textit{prefer} algorithmic to human advice in predictive domains, exhibiting what they term ``algorithm appreciation.'' A unidimensional model treats these as contradictory findings requiring reconciliation through moderator variables: perhaps aversion appears in some contexts and appreciation in others, depending on task type, outcome observability, or individual differences.

We propose that the contradiction is an artifact of dimensional collapse. Algorithm appreciation and algorithm aversion are not opposing ends of a single trust continuum; they are manifestations of different trust \textit{modes} operating under different conditions. Investors who evaluate algorithmic performance against benchmarks are engaged in a qualitatively different trust relationship from investors who grant an algorithmic advisor discretionary authority over their portfolio. The former --- which we term \textit{evaluative trust} --- is inherently conditional on continued satisfactory performance and produces rapid exit when performance falters. The latter --- which we term \textit{delegative trust} --- involves a transfer of judgment that is not reducible to performance monitoring and produces a qualitatively different set of responses to adverse events, including both greater resilience under some conditions and a distinctive \textit{betrayal} reaction under others.

The distinction between evaluative and delegative trust is not merely taxonomic. It generates predictions that the unidimensional model cannot. If trust varies only in degree, then higher-trust investors should uniformly exhibit greater resilience after adverse performance: they trusted more, so they tolerate more. But practitioners and emerging empirical evidence suggest a more complex pattern. Some investors with strong pre-shock trust relationships are among the \textit{most} reactive to performance failures, displaying emotional responses --- a sense of violated expectations, of misplaced faith --- that exceed what a rational updating model would predict. The unidimensional framework has no vocabulary for this phenomenon. The evaluative/delegative framework does: these investors extended \textit{delegative} trust, and its violation produces not the disappointment of a failed prediction but the betrayal of a breached relationship.

% --- Move 2: The Distinction ---

\subsection{The Distinction}
\label{subsec:distinction}

The evaluative/delegative framework draws on a substantial body of philosophical work on the nature of trust. \citet{CarterSpencerSimion2020} demonstrated the relevance of philosophical trust concepts to financial contexts; the present paper provides the first empirical operationalization of these concepts with original data.

\subsubsection{Evaluative Trust}
\label{subsubsec:evaluative_trust}

Evaluative trust is the mode of trust that operates through assessment, verification, and conditional continuation. An investor operating in evaluative trust mode monitors portfolio performance, compares returns against relevant benchmarks, examines the rationale for specific allocation decisions, and maintains the advisory relationship conditional on continued evidence of competent management. The relationship is governed by evidence: when evidence supports the trustee's competence, the relationship continues; when evidence no longer supports it, the investor withdraws.

This mode of trust maps to well-established constructs in both philosophy and organizational behavior. \citet{Hardin2002} models trust as ``encapsulated interest'': A trusts B because A believes that B's interests encapsulate A's interests. That is, B will perform well for A because doing so serves B's own goals. The trust is grounded in a prediction about behavior derived from an assessment of incentive alignment, and it is withdrawn when the incentive structure no longer supports the prediction. \citet{McAllister1995} identifies a cognition-based dimension of interpersonal trust in organizations, grounded in evidence of the trustee's competence and reliability, which he distinguishes from an affect-based dimension grounded in emotional bonds and mutual care. Evaluative trust in financial contexts aligns closely with \citeauthor{McAllister1995}'s cognition-based construct: it is rational, evidence-conditional, and does not require emotional attachment to the trustee.

Evaluative trust is not inherently shallow or deficient. Sophisticated evaluative trust --- in which an investor develops nuanced frameworks for assessing risk-adjusted performance, understands factor exposures, and monitors for style drift --- represents a high level of financial engagement. A client who retains an advisor because they have carefully evaluated the advisor's strategy and believe it sound is exercising evaluative trust. The key feature is not the depth of engagement but the \textit{conditionality}: the relationship is maintained because and insofar as the evidence supports it.

The behavioral signatures of evaluative trust follow directly from its structure. Because evaluative trust is conditional on evidence, its withdrawal is proportional to the strength of disconfirming evidence and takes the form of \textit{disappointment} --- a cognitive assessment that expectations were not met. The investor's response to a performance shock is to seek information (Was the underperformance systematic or idiosyncratic? How did comparable strategies perform?), reassess the relationship, and either revise expectations or exit. The emotional register is analytical: the investor feels let down, not betrayed.

\subsubsection{Delegative Trust}
\label{subsubsec:delegative_trust}

Delegative trust involves a qualitatively different relationship. An investor operating in delegative trust mode has granted the advisor --- whether human or algorithmic --- discretionary authority over decisions that the investor cannot or chooses not to fully evaluate. The investor has, in a meaningful sense, \textit{transferred judgment}: not merely delegated execution of a strategy the investor has approved, but ceded the authority to determine what strategy is appropriate, to adjust it as conditions change, and to exercise professional judgment about matters the investor does not independently assess.

This mode of trust draws on several philosophical accounts that have not previously been applied to financial relationships. \citet{Baier1986} argues that trust inherently involves accepted vulnerability: to trust is to be in a position where the trusted party can harm you, and to accept that position rather than eliminate it through control or monitoring. \citet{Nickel2012}, analyzing trust in medical contexts, characterizes it as entrusting someone with discretionary authority --- authority to make decisions using their own judgment, not merely to execute pre-specified instructions. \citet{Hawley2014} adds that trust involves a \textit{commitment}: to trust someone is to rely on their commitment to do what they are relied upon to do, and to feel betrayal when that commitment is violated. The trusted party is not merely predicted to perform; they are counted on to fulfill an obligation.

\citet{Holton1994} provides the critical behavioral distinction, drawing on \citeauthor{Strawson1962}'s~(\citeyear{Strawson1962}) framework of reactive attitudes. He argues that what differentiates trust from mere reliance is the \textit{reactive attitude} of the trustor upon violation. When I rely on my alarm clock and it fails, I am annoyed --- but I do not feel betrayed, because the clock made no commitment and I invested no vulnerability. When I trust my financial advisor and she makes unauthorized trades, I feel betrayed --- because I granted her discretionary authority, accepted the vulnerability that comes with not monitoring every decision, and she violated the commitment that this acceptance implied. The presence or absence of betrayal as a reactive attitude is, for Holton, the diagnostic marker of whether a relationship constitutes trust or mere reliance.

Applied to financial relationships, delegative trust manifests as acceptance of opacity (the investor does not insist on understanding every decision), tolerance of short-term uncertainty (the investor does not react to every quarterly fluctuation), and a distinctive \textit{scope} of the relationship that extends beyond narrow performance metrics. The finding by \citet{RossiUtkus2020} that advisory clients cite ``peace of mind'' and ``delegation'' as primary reasons for hiring advisors points toward delegative trust, though the existing literature does not isolate it from evaluative trust.

The behavioral signatures of delegative trust differ sharply from those of evaluative trust, particularly in response to adverse events. Because the investor has accepted vulnerability and granted discretionary authority, violation produces not proportional disappointment but \textit{betrayal} --- a response whose intensity reflects the depth of the trust relationship rather than the magnitude of the performance shortfall. The betrayal response is qualitatively distinct: the investor feels not merely that their predictions were wrong, but that their vulnerability was exploited, that the commitment they relied upon was hollow. This is the phenomenon that the unidimensional model cannot capture --- the investor who was \textit{most} trusting in the delegative sense may produce the \textit{most} intense negative reaction, because they had the most to lose.

A further distinction concerns what happens \textit{between} adverse events. An evaluative trustor monitors continuously; their trust is maintained through ongoing verification. A delegative trustor monitors less --- not from inattention or ignorance, but from a \textit{choice} to refrain from verification that would convert the relationship from delegation to surveillance. This distinction generates a testable prediction: platforms that encourage frequent performance monitoring may inadvertently shift users from delegative to evaluative trust mode, increasing the likelihood of exit after subsequent performance shocks. \citet{Pettit1995} argues that introducing checks and surveillance removes the opportunity to manifest trusting reliance, thereby foreclosing the trust relationship itself; the monitoring that was meant to ensure compliance prevents trust from forming in the first place. The investor who checks their portfolio daily is, through the act of checking, operating in evaluative mode --- regardless of how the relationship was initially framed.

Delegative trust also has a distinctive failure mode that is not simply the inverse of its formation. \citet{hardin1993street} observes that learned extreme distrust is self-reinforcing: it forecloses the encounters that might update the prior. A person who has developed generalized distrust of a category of agent avoids interactions with members of that category, ensuring they never accumulate the experience that could revise the disposition. Applied to our framework, this produces a phenomenon we term \textit{delegative foreclosure}: the closure of an entire trust mode toward a category of trustee, not the downward revision of a trust estimate within that mode. A consumer who has had a negative experience with one algorithmic advisor --- or who has absorbed sufficiently negative narratives about the category --- does not merely lower their evaluative assessment of algorithmic advice. They foreclose the \textit{possibility} of extending delegative trust to any algorithmic system. The evaluative channel remains open: they can still be shown performance data, benchmark comparisons, and track records, and they may revise their evaluative trust accordingly. But the delegative channel is closed. No amount of performance evidence is likely to reopen it, because performance evidence addresses the wrong dimension.

Delegative foreclosure is categorically different from algorithm aversion as the existing literature describes it. Dietvorst-style aversion is an evaluative phenomenon: the user observed poor performance and updated their competence assessment downward. Foreclosure is a structural change in the trustor's available modes of relating to algorithmic trustees. This distinction matters practically because the interventions appropriate to each are different. Evaluative aversion can be addressed through better performance data, transparency, and track records. Foreclosure requires a different kind of repair --- the reconstruction of conditions under which the trustor can experience vulnerability as safe, which is a relational project, not an informational one.

\subsubsection{Formal Characterization}
\label{subsubsec:formal}

The preceding accounts can be stated with greater precision. Let investor $i$'s trust relationship with advisor $j$ in financial domain $D$ be characterized by an ordered pair $(E, \delta)$, where each component captures a distinct dimension of the relationship.

\textbf{Definition 1 (Evaluative Trust).} $E(i, j, D) \in [0, 1]$ is investor $i$'s evidence-conditional assessment of advisor $j$'s competence in domain $D$. Let $s_t$ denote a performance signal received at time $t$ (e.g., portfolio returns relative to a benchmark). Evaluative trust updates in response to performance information:
\[
E_t = f(E_{t-1}, s_t), \quad \text{where } \frac{\partial f}{\partial s} > 0.
\]
The function $f$ is monotonically increasing in the signal: positive performance signals raise evaluative trust; negative signals lower it. When $E_t$ falls below investor $i$'s continuation threshold $E^*$, the investor exits. The reactive attitude upon exit is \textit{disappointment}, proportional to the magnitude of the shortfall $(E^* - E_t)$.

\textbf{Definition 2 (Delegative Trust).} $\delta(i, j, D) \in [0, 1]$ is the degree to which investor $i$ has granted advisor $j$ discretionary authority over decisions in $D$ that $i$ does not independently verify. Unlike evaluative trust, delegative trust is not updated by performance signals alone:
\[
\frac{\partial \delta}{\partial s_t} \approx 0 \quad \text{for ordinary performance signals.}
\]
Delegative trust responds instead to \textit{relational signals} --- evidence about the commitment structure itself. Let $r_t$ denote a relational signal (e.g., discovery of a fiduciary breach, unauthorized exercise of discretion, institutional restructuring that alters the authority relationship). Then:
\[
\delta_t = h(\delta_{t-1}, r_t), \quad \text{where } \frac{\partial h}{\partial r} \text{ is large and negative for commitment violations.}
\]
The reactive attitude upon violation is \textit{betrayal}, whose intensity reflects the prior level of delegative trust: a high-$\delta$ investor who experiences a relational violation produces a stronger negative response than a high-$E$ investor who experiences an equivalent performance shock, because what has been violated is not a prediction but a grant of discretionary authority.

These definitions generate three claims that the subsequent empirical analyses test.

\textbf{Claim 1 (Dimensional Separability).} $E$ and $\delta$ respond to different signal classes. Across contexts that vary the salience of performance information versus relational commitment, $E$ and $\delta$ should vary independently. Formally: a manipulation that increases the salience of performance monitoring (e.g., benchmark framing) should raise $E$ without proportionally raising $\delta$, while a manipulation that emphasizes fiduciary commitment should raise $\delta$ without proportionally raising $E$. A unidimensional construct would require both to move together.

\textbf{Claim 2 (Asymmetric Post-Shock Response).} After an adverse event $s_t < \bar{s}$:
\begin{quote}
High-$E$ investors: $\Pr(\text{exit}) = g(\bar{s} - s_t, E_{t-1})$ --- exit probability increasing in shortfall magnitude.

High-$\delta$ investors: $\Pr(\text{exit})$ is low when $s_t$ carries only performance information, because $\partial \delta / \partial s_t \approx 0$. But if $s_t$ also conveys relational information $r_t$ (the advisor failed to exercise genuine judgment, the institutional commitment was hollow), then $\Pr(\text{exit})$ can \textit{exceed} the evaluative case, accompanied by betrayal rather than disappointment.
\end{quote}
This generates what we term the betrayal premium prediction --- related to the pattern documented by \citet{kormylo2026}, but reframed through the evaluative/delegative distinction: for investors in the delegative mode, the gap between betrayal intensity and disappointment intensity should be positive after an adverse event, while for evaluative investors it should be near zero or negative.

\textbf{Claim 3 (The Transparency Paradox).} Let $\tau$ represent the degree of transparency or auditability in the advisory relationship. Increasing $\tau$ shifts the trust relationship along different dimensions:
\begin{align*}
\frac{\partial E}{\partial \tau} &> 0 \quad \text{(transparency makes performance verifiable, supporting evaluative trust)} \\
\frac{\partial \delta}{\partial \tau} &< 0 \quad \text{(salient, frequent monitoring erodes the discretionary space in which delegative trust operates)}
\end{align*}
This is the mechanism underlying \citeauthor{ONeill2002}'s~(\citeyear{ONeill2002}) diagnosis that audit culture undermines trust, formalized through the two-dimensional framework. A unidimensional model predicts that transparency uniformly increases trust; the two-dimensional model predicts that it increases one component while decreasing the other --- with net effects on investor behavior that depend on which mode dominates the relationship.

\subsubsection{The Relationship Between the Two Modes}
\label{subsubsec:relationship}

The formal characterization makes explicit that evaluative and delegative trust are distinct dimensions, not opposite poles of a single continuum. In principle, an investor can exhibit high levels of both: carefully evaluating an advisor's track record (high $E$) while simultaneously granting discretionary authority over tactical decisions within an approved strategic framework (high $\delta$). An investor can also exhibit low levels of both --- maintaining a financial relationship through inertia without either evaluating performance or granting meaningful discretionary authority.

However, the two modes are not fully independent in practice. In constrained communicative contexts --- including the complaint narratives analyzed in this study --- the two modes compete for expression. A narrative dominated by relational language (betrayal, violated expectations, exploited vulnerability) leaves less space for analytical language (benchmark comparison, performance metrics, cost assessment), and vice versa. This produces an inverse correlation in observed expression that reflects the communicative constraint rather than a structural opposition between the constructs themselves. The critical test of dimensional separability is not the within-context correlation but the \textit{cross-context asymmetry} predicted by Claim~1: whether the two dimensions vary independently across different conditions. As we demonstrate in the empirical analysis, delegative trust expression varies substantially across financial product categories while evaluative trust expression remains essentially constant --- a pattern that a single bipolar dimension cannot produce.

The distinction between evaluative and delegative trust also clarifies a long-standing measurement problem. \citeauthor{McAllister1995}'s~(\citeyear{McAllister1995}) affect-based trust items --- developed for interpersonal workplace relationships --- fail when applied to algorithmic trustees because they operationalize the relational dimension of trust through interpersonal warmth (``We have a sharing relationship''; ``I can talk freely to this individual about difficulties''). This failure has been interpreted as evidence that algorithms cannot receive trust beyond the competence-based dimension.%
\footnote{See, e.g., the discussion in \citet{LeeSee2004} and \citet{Muir1994} on the limits of trust in automated systems. \VERIFYSAGE{Confirm these are the best citations for the claim that algorithms cannot receive trust beyond competence-based trust; add or replace as needed.}} We argue that the conclusion is premature. The failure is one of \textit{operationalization}, not of \textit{construct}: delegative trust with algorithmic systems operates not through affective bonds but through institutional features --- fiduciary framing, regulatory protections, discretionary authority structures --- that create the conditions for vulnerability acceptance and judgment transfer without requiring interpersonal warmth. In terms of the formal characterization, institutional features can raise $\delta$ even when affective channels are absent, provided they create credible commitment structures. The Prolific experiment reported in this paper tests this claim directly.

% --- Move 3: The Agency Theory Bridge ---

\subsection{The Agency Theory Bridge}
\label{subsec:agency_bridge}

The evaluative/delegative distinction maps onto a foundational tension in financial economics: the spectrum between complete and incomplete contracting. Under \citeauthor{JensenMeckling1976}'s~(\citeyear{JensenMeckling1976}) framework, the principal--agent relationship is governed by a contract that specifies observable performance metrics, monitoring mechanisms, and enforcement provisions. The principal's problem is one of verification --- ensuring the agent acts as contracted --- and the solution is better information. This is evaluative trust rendered in the language of corporate finance: the investor trusts the advisor because and insofar as performance is observable, contractible, and enforceable.

But real advisory relationships are not complete contracts. \citet{GrossmanHart1986} demonstrate that when contingencies cannot be fully specified in advance --- when the space of possible market states exceeds the contracting parties' ability to enumerate them --- residual control rights must be allocated. Someone must decide what to do in circumstances the contract does not address. Delegative trust, in this framing, is the investor's acceptance of residual control rights held by the advisor: the authority to exercise judgment over decisions that were not and could not have been contracted upon. The investor does not merely delegate execution of a pre-specified strategy; they cede the right to determine what strategy is appropriate as conditions change. The trust inheres not in the ability to verify outcomes but in the willingness to grant discretion over uncontracted states of the world.

This mapping carries a methodological implication. The trust game introduced by \citet{BergDickhautMcCabe1995} --- the dominant experimental paradigm for measuring trust in economics --- is structurally a complete-contracting environment. The payoff matrix is known, the interaction is typically one-shot, the counterparty is anonymous, and the decision reduces to a single monetary transfer with a calculable expected return. Multi-round and repeated-interaction variants exist \citep{JohnsonMislin2011}, but even these retain the core feature of observable, calculable payoffs --- the trustor always knows the outcome of each round and can condition future behavior on it. These are precisely the conditions under which evaluative trust operates: the trustor assesses the probability of reciprocation, calibrates the amount sent accordingly, and withdraws from the interaction once the outcome is observed. There is no residual control, no unspecified contingency, no discretionary authority to grant. If the finance literature validates its trust measures against trust game behavior --- and \citeauthor{JohnsonMislin2011}'s~(\citeyear{JohnsonMislin2011}) meta-analysis documents extensive use of the paradigm across disciplines --- then the empirical foundation may be systematically measuring evaluative trust while the advisory relationships it seeks to explain are substantially delegative.

% --- Move 4: The "Peace of Mind" Reinterpretation ---

\subsection{The ``Peace of Mind'' Reinterpretation}
\label{subsec:peace_of_mind}

The evaluative/delegative framework's explanatory power is most clearly demonstrated against the most influential model of trust in financial advice. \citet{GennaioliShleiferVishny2015} propose that financial advisors function as ``money doctors'' whose primary service is the provision of trust. In their model, trust reduces the client's \textit{perceived} variance of returns:\VERIFYSAGE{Confirm whether GSV (2015) say ``perceived variance of returns'' or ``variance of expected returns.''} the advisor does not necessarily improve expected outcomes, but the client experiences the portfolio as less risky because a trusted professional is managing it. Clients pay advisory fees not for alpha but for ``peace of mind.'' The model is elegant, empirically productive, and --- we argue --- ambiguous on a dimension that determines its behavioral predictions.

The ambiguity is this: what kind of trust reduces perceived variance? Under an evaluative interpretation, the advisor provides peace of mind through optimization and legibility. The client monitors less frequently not because they have transferred judgment but because the advisor's demonstrated competence has lowered the expected incidence of events worth monitoring. The client's comfort is evidence-conditional: it persists because performance has been satisfactory and would dissolve if performance deteriorated. Under a delegative interpretation, peace of mind arises from the act of judgment transfer itself. The client experiences reduced variance not because outcomes are objectively smoother but because the burden of decision-making --- of choosing, second-guessing, and bearing responsibility for allocation choices --- has been lifted. The comfort comes from the relationship, not from verification of the relationship's outputs.

These interpretations generate sharply different predictions about what happens after a performance shock. The evaluative peace-of-mind client experiences the shock as a variance event that exceeds prior expectations. Their response is to reassess: Was the advisor's strategy flawed? Did the underperformance reflect systematic failure or market conditions? The emotional register is disappointment calibrated to the magnitude of the shortfall relative to expectations. If the evidence warrants it, they exit; if it does not, they revise expectations and continue. The delegative peace-of-mind client experiences the shock differently. If the advisory relationship remains intact --- if the advisor communicates, explains the approach going forward, and continues to exercise the discretionary authority the client granted --- the client's response may be patience. The advisor is navigating difficulty on the client's behalf, which is precisely what they were trusted to do. But if the shock reveals that the delegative relationship was hollow --- that the advisor was not exercising genuine judgment, or that the institutional structure does not support the discretionary authority the client believed they had granted --- the response is betrayal, not disappointment. The intensity of the reaction reflects the depth of the prior vulnerability, not the magnitude of the financial loss.

Survey evidence suggests the delegative interpretation better describes advisory clients' self-reported experience. \citet{RossiUtkus2020} find that clients cite ``delegation'' and ``peace of mind'' as primary reasons for retaining advisors --- language that maps more naturally to judgment transfer than to variance assessment. But the existing literature does not isolate the two mechanisms. The Money Doctors model is agnostic about the \textit{kind} of trust that does the variance-reducing work, and its predictions are correspondingly silent about the qualitative differences in post-shock behavior that the evaluative/\allowbreak delegative distinction reveals. Our framework does not replace the \citeauthor{GennaioliShleiferVishny2015} model; it disaggregates a parameter the model treats as unitary.

% --- Move 5: Reconciling Algorithm Appreciation and Aversion ---

\subsection{Reconciling Algorithm Appreciation and Aversion}
\label{subsec:reconciling}

Three findings in the algorithmic decision-making literature appear to be in tension. \citet{LoggMinsonMoore2019} find that individuals systematically prefer algorithmic to human advice in predictive domains, exhibiting ``algorithm appreciation.'' \citet{DietvorstSimmonsMassey2015} find that individuals who observe an algorithm err subsequently refuse to use it, even when it outperforms human forecasters --- ``algorithm aversion.'' And \citet{BigmanGray2018} find that individuals resist algorithmic involvement in decisions with moral dimensions, even when the algorithm's judgments are equivalent to or better than human ones. The literature has treated these as context-dependent findings requiring reconciliation through moderator variables\VERIFYSAGE{Add citations for the moderator-variable reconciliation literature on algorithm aversion vs.\ appreciation; flag for author if none found.}: appreciation prevails in some task environments, aversion in others, depending on outcome observability, task subjectivity, and the moral valence of the decision.

The evaluative/delegative framework offers a more parsimonious resolution. As documented in the existing literature, algorithm appreciation is evaluative trust operating as expected: when an algorithm demonstrates predictive competence in a domain where performance is observable and quantifiable, evaluative trust increases. The algorithm is assessed, verified, and preferred precisely because it outperforms. Algorithm aversion after errors is evaluative trust failing as expected: the same performance-conditional relationship that produced appreciation now produces withdrawal when a negative signal arrives. The trust was never resilient to error because it was never grounded in anything beyond continued evidence of competence. Both findings, as the literature describes them, document evaluative trust under different performance conditions, not opposing orientations toward algorithms. Our framework leaves open the possibility that delegative trust toward algorithms can also form --- as the Prolific experiment tests directly --- but the existing appreciation/aversion literature has not studied this dimension.

Resistance to algorithmic moral decisions is a different phenomenon entirely. Moral decisions require precisely the exercise of discretionary judgment that defines delegative trust: they cannot be reduced to optimization over a known objective function, they require weighing incommensurable values, and they implicate the decision-maker's character --- their integrity, not merely their competence. \citeauthor{BigmanGray2018}'s finding may reflect, at least in part, resistance to extending \textit{delegative} trust to algorithms, rather than a generalized aversion. Respondents who willingly delegate predictive tasks to an algorithm (evaluative trust) balk at delegating moral judgment to the same system (delegative trust). The trust mode, not the trustee, has changed.

The concept of delegative foreclosure introduced above explains a further puzzle in this literature: the persistence of algorithm aversion even when users are presented with compelling evidence of algorithmic superiority. If the resistance were purely evaluative --- a low competence assessment --- then strong performance data should resolve it. \citet{DietvorstSimmonsMassey2015} find that it does not. Our framework identifies the mechanism: for users who have foreclosed delegative trust toward algorithms, performance evidence addresses the evaluative dimension while the operative barrier is structural --- the delegative channel is closed, and no volume of evidence within the evaluative channel is likely to reopen it. The appropriate intervention is not more data but the reconstruction of conditions under which delegation feels safe.

\citeauthor{DietvorstSimmonsMassey2018}'s~(\citeyear{DietvorstSimmonsMassey2018}) follow-up study provides further support. They find that even trivial modification ability --- the capacity to adjust an algorithm's output by a small amount before it takes effect --- substantially reduces algorithm aversion. This result is puzzling under a pure competence account: a minor adjustment does not meaningfully improve the algorithm's performance, so why should it change the user's willingness to rely on it? Under the evaluative/delegative framework, the answer is immediate. Modification grants the user a form of discretionary authority over the final output. In \citeauthor{Nickel2012}'s~(\citeyear{Nickel2012}) terms, the user is no longer a passive recipient of algorithmic decisions but a participant who exercises judgment --- however minimally --- over the outcome. This converts the relationship from pure evaluation (observing and assessing the algorithm's performance) to partial delegation (sharing authority over the final decision). The trust changes in kind, not merely in degree, and the behavioral consequences follow: a user who has exercised even token authority over a decision is less likely to experience the algorithm's error as a failure imposed upon them and more likely to experience it as a joint outcome for which they bear partial responsibility. \VERIFYSAGE{Confirm that Dietvorst et al.\ (2018) is the correct citation for the modification-ability follow-up, and whether the effect was characterized as ``trivial'' or ``small'' modification.}

% --- Move 6: The Transparency Paradox and Hypotheses ---
% TODO: Review this section for repetition with Discussion §6.4 (Delegative Foreclosure
% and Design Implications). Consider tightening the O'Neill material here and expanding
% the connection to the trust-in-finance literature (e.g., Guiso/Sapienza/Zingales,
% Gennaioli/Shleifer/Vishny) to differentiate the theoretical prediction from the
% empirical discussion. —CT revision note 2026-03-03

\subsection{The Transparency Paradox and Hypotheses}
\label{subsec:transparency}

The evaluative/delegative framework generates a counterintuitive prediction about a central policy question in financial regulation: the role of transparency. \citet{ONeill2002, ONeill2014} argues that the contemporary culture of transparency --- the proliferation of disclosure mandates, performance metrics, audit requirements, and accountability mechanisms --- has failed to restore public trust in institutions despite being designed expressly for that purpose. Her diagnosis is that transparency and audit culture cannot solve a trust crisis, and may in fact deepen it. The mechanisms of accountability, intended to make institutions more trustworthy, have instead produced ``a culture of suspicion'' in which institutional actors optimize for measurability rather than the underlying goods they are meant to deliver.

\citeauthor{ONeill2002}'s diagnosis is acute but theoretically incomplete. She identifies the failure without fully specifying the mechanism by which transparency undermines trust. The evaluative/delegative framework supplies this mechanism. Transparency builds evaluative trust: it makes the trustee's actions observable, verifiable, and assessable against explicit criteria. Disclosure mandates, benchmark reporting, and performance analytics all serve to make the advisory relationship more legible to the investor. But transparency simultaneously erodes the conditions under which delegative trust operates. Delegative trust requires discretionary space --- the authority to exercise judgment over matters that the principal has chosen not to fully monitor. When every decision must be disclosed, justified, and subjected to audit, the discretionary space collapses. The advisor no longer holds residual control rights; they execute under surveillance. And an investor who verifies every decision is, by definition, operating in evaluative mode --- regardless of how the relationship was initially framed or intended.

Applied to algorithmic finance, this analysis yields a specific prediction: explainability mandates for robo-advisors --- regulatory requirements that algorithms disclose the reasoning behind their investment decisions --- may paradoxically \textit{increase} algorithm aversion rather than reduce it. By making algorithmic decisions transparent and auditable, regulators shift users from a delegative trust equilibrium (in which they accept the algorithm's discretionary authority and tolerate short-term uncertainty) to an evaluative trust equilibrium (in which they assess each decision against their own expectations and benchmarks). The shift increases the user's sensitivity to short-term performance fluctuations and lowers the threshold for exit --- precisely the opposite of the regulatory intent. Mandatory performance disclosure may induce monitoring behavior that converts the long-term commitment benefits documented by \citet{GennaioliShleiferVishny2015} into the short-term reactivity documented by \citet{DietvorstSimmonsMassey2015}.

\citet{CarterSpencerSimion2020} identified the need for philosophical trust theory in financial contexts, applying \citeauthor{Baier1986}'s vulnerability account and \citeauthor{Hawley2014}'s commitment framework to the 2008 banking crisis. Their analysis demonstrated that reactive attitudes --- betrayal, resentment, a sense of violated expectations --- distinguish trust failures from mere reliance failures in financial relationships. But their contribution remained at the level of diagnosis: they showed that philosophical trust concepts apply to finance without operationalizing the distinction, without testing whether institutional framing shifts investors between trust modes, and without extending the framework to the algorithmic context that now dominates retail financial advice. The present paper answers all three.

\subsubsection{Hypotheses}
\label{subsubsec:hypotheses}

The theoretical framework developed across the preceding sections generates five pre-registered predictions, tested through a between-subjects experiment in which participants are randomly assigned to one of three conditions --- fiduciary framing (Condition A), performance framing (Condition B), or a neutral control (Condition C) --- before encountering an identical adverse performance event.

\textbf{H1 (Trust Mode Separation).} Fiduciary framing produces higher delegative trust than performance framing; performance framing produces higher evaluative trust than fiduciary framing. If evaluative and delegative trust are distinct dimensions responsive to different institutional features, then framing that emphasizes discretionary authority and relational commitment (Condition A) should elevate delegative trust, while framing that emphasizes verifiable performance and benchmark comparison (Condition B) should elevate evaluative trust.

\textbf{H2 (Betrayal Premium).} The difference between betrayal and disappointment responses after the adverse event is larger under fiduciary framing than under performance framing. Delegative trust violation produces betrayal~\citep{Holton1994}; evaluative trust failure produces disappointment. If Condition A successfully induces delegative trust, the betrayal response should be elevated relative to disappointment --- a pattern we term the ``betrayal premium.''

\textbf{H3 (Algorithm Aversion Mechanism).} The tendency to switch from algorithmic to human advice after the adverse event is predicted by evaluative trust scores but not by delegative trust scores. Algorithm aversion, in the Dietvorst sense, is an evaluative trust phenomenon: performance-conditional trust withdrawn after a negative signal. Delegative trust, grounded in accepted vulnerability and granted discretion rather than in performance verification, should not predict the algorithm-to-human switch.

\textbf{H4 (Delegative Resilience).} Higher delegative trust scores predict lower exit-seeking behavior after the adverse event, controlling for evaluative trust. Delegative trust tolerates short-term uncertainty because the trustor has accepted the trustee's exercise of discretionary authority through adverse conditions --- the advisor is navigating difficulty, which is what they were trusted to do.

\textbf{H5 (Financial Literacy Interaction).} The treatment effect of fiduciary versus performance framing on delegative trust is concentrated among high-literacy respondents. Low-literacy respondents default to elevated delegative scores across conditions --- not from trust, but from the absence of evaluative capacity. High-literacy respondents \textit{can} evaluate but, under fiduciary framing, are given institutional permission not to. The interaction distinguishes genuine delegation --- a choice made by those with the capacity to do otherwise --- from delegation by default.
